% UHH Assignment essay starter — uhhassignment.cls (LPPL-1.3c)
% The class is supplied by TeX Live; this document is an original example.
\documentclass[shorttitle,repeatauthor]{uhhassignment}

\usepackage[T1]{fontenc}
\usepackage[utf8]{inputenc}
\usepackage{amsmath,amssymb}
\usepackage{booktabs}
\usepackage{enumitem}
\usepackage{microtype}

\university[Northbridge University]{Northbridge University}
\course[Environmental Policy]{Environmental Policy and Evidence}
\teacher{Dr. Lina Okafor}
\group{B}
\sheetno{3}
\author[Jordan Blake]{Jordan Blake \id{j.blake@example.edu}}
\title{\textsc{Evidence, uncertainty and public decisions}}
\date{18 September 2026}

\begin{document}
\maketitle

\task{Framing the decision}
Public policy arguments often present a forecast as though it were a
measurement.  This essay distinguishes the two.  A measurement describes
an observed quantity; a forecast combines observations with a model and
assumptions about what will happen next.  That difference is not a reason
to abandon quantitative evidence.  It is a reason to expose the assumptions
that make the evidence useful.

\subtask{A practical distinction}
Suppose a city records the daily number of bicycle trips before and after a
protected lane opens.  The observed change is not, by itself, the effect of
the lane: weather, fuel prices and a nearby road closure may have changed at
the same time.  A defensible analysis therefore states a comparison group,
the time window, and the mechanism by which the intervention is expected to
change behaviour.

\task{A small model}
Let $Y_1$ and $Y_0$ be the number of bicycle trips on a representative day
with and without the lane.  The policy effect is the counterfactual

\[
  \tau = \mathbb{E}[Y_1-Y_0].
\]

Only one of the two potential outcomes can be observed for a given day.
That missing counterfactual is why a comparison design matters.  If a
nearby untreated route has a stable trend, a difference-in-differences
estimate can be written as

\[
 \widehat{\tau}= (\overline{Y}_{\mathrm{treated,after}}
 -\overline{Y}_{\mathrm{treated,before}})
 -(\overline{Y}_{\mathrm{control,after}}
 -\overline{Y}_{\mathrm{control,before}}).
\]

The estimate is informative only if the two routes would have followed
parallel trends in the absence of the intervention.  That assumption is a
claim about the design, not a number produced by the formula.

\subtask{Interpretation}
\begin{enumerate}[label=\alph*),leftmargin=2em]
  \item State the target population before calculating an effect.
  \item Report an uncertainty interval and the number of observations.
  \item Test whether the result changes under a plausible alternative
        comparison route.
\end{enumerate}

\task{Reflection}
The strongest conclusion available from this small study is conditional:
if the comparison route captures the city-wide time trend and no other
route-specific change occurred, the protected lane is associated with an
increase in cycling.  Calling that association a causal effect without
checking those conditions would turn a useful model into overconfident
storytelling.  The assignment's reproducible layout makes those conditions
visible alongside the argument, which is the point of a good hand-in.

\end{document}
